% Appendix: the Lean formalisation. Input from main.tex after \appendix.
\section{The Lean Formalisation}\label{app:lean}

This appendix describes the mechanised model behind \S\ref{sec:security}
in enough detail that a reader can judge what has been proved, under
what assumptions, and what has not. The development is about 1{,}300
lines of Lean~4~\cite{lean4} in seven modules and depends on nothing but the
core library: no Mathlib, no external tactics. It builds in seconds and
prints an axiom audit for every headline theorem; each depends only on
\lean{propext}, \lean{Quot.sound} and, in one case, \lean{Classical.choice},
and none on \lean{sorryAx}. The code is in the artefact~\cite{layertreerepo}.

\subsection{Design of the Model}

The model states the \emph{rules} of \S\ref{sec:transitions} as
relations over the objects the scripts manipulate: nodes, paths, amounts.
It does not model tapscript bytecode, signatures, keys or lock times.
Authorisation is carried as an opaque proposition, and the input and
output amounts the scripts read with \op{TX} are parameters. Three
choices keep the proofs short.

\paragraph{Depth as a type index.} A tree of depth $d$ is
\lean{Tree d}, a perfect binary tree whose tips are \lean{Option Leaf}.
Positions are \lean{Pos d h}, paths from depth $d$ down to a subtree of
height $h$, so that \lean{Pos d 0} addresses a leaf and \lean{Pos d h}
addresses the slot a \tr{JOIN}$_h$ fills. Indexing by both endpoints
avoids all arithmetic on depths: the recursions on \lean{getSub},
\lean{setSub} and \lean{pathSub} pattern-match on the position and the
tree simultaneously, and the one impossible case (a \lean{cons} step below
height~$0$) is dismissed by \lean{Pos.height\_le}.

\begin{lstlisting}
inductive Tree : Nat → Type
  | tip : Option Leaf → Tree 0
  | bin {d : Nat} : Tree d → Tree d → Tree (d + 1)

inductive Pos : Nat → Nat → Type
  | nil {h : Nat} : Pos h h
  | cons {d h : Nat} : Bool → Pos d h → Pos (d + 1) h

def parent (a b : Node H) : Node H :=
  ⟨M.hNode a.hash a.sum b.hash b.sum, a.sum + b.sum⟩

def combine (n : Node H) (x : Bool × Node H) : Node H :=
  if x.1 then parent x.2 n else parent n x.2

def fold (n : Node H) : List (Bool × Node H) → Node H
  | [] => n
  | x :: rest => combine (fold n rest) x
\end{lstlisting}

\paragraph{Root-first paths.} The model orders a sibling path from the
root down, so \lean{fold} is a structural recursion on the list and the
soundness proof is a plain induction. The witness of \S\ref{sec:scripts}
supplies the same pairs leaf-first, because the script consumes the
deepest level first; nothing in the theorems depends on the order.

\paragraph{Hashes as parameters.} Everything is parametric in a class
\lean{HashModel} whose fields are the two hash functions, the zero digest,
and the assumptions the argument uses:

\begin{lstlisting}
class HashModel (Leaf : Type) (H : outParam Type) where
  hLeaf : Leaf → H
  hNode : H → Nat → H → Nat → H
  empty0 : H
  hLeaf_inj : ∀ {l₁ l₂}, hLeaf l₁ = hLeaf l₂ → l₁ = l₂
  hNode_inj : ∀ {a₁ b₁ a₂ b₂ sa₁ sb₁ sa₂ sb₂},
    hNode a₁ sa₁ b₁ sb₁ = hNode a₂ sa₂ b₂ sb₂ →
      a₁ = a₂ ∧ sa₁ = sa₂ ∧ b₁ = b₂ ∧ sb₁ = sb₂
  hNode_ne_hLeaf : ∀ a sa b sb l, hNode a sa b sb ≠ hLeaf l
  hNode_ne_empty : ∀ a sa b sb, hNode a sa b sb ≠ empty0
  hLeaf_ne_empty : ∀ l, hLeaf l ≠ empty0
\end{lstlisting}

Injectivity is collision resistance idealised; the three inequalities are
the tag separation that distinct BIP-340 tags provide. The encodings
module discharges the serialisation half of each injectivity field
(Lemma~\ref{lem:enc}): given SHA-256 collision resistance,
\lean{serializeLeaf\_inj} and \lean{serializeNode\_inj} are what remains.
Nothing quantifies the adversary's advantage; a proof under injectivity
transfers to any adversary that cannot exhibit a collision.

\subsection{Theorem Statements}

The statements below are verbatim from the development. \lean{t.root} is
the root node of a tree, \lean{t.getSub p} the subtree at position
\lean{p}, \lean{t.setSub p u} the tree with that subtree replaced by
\lean{u}, \lean{t.pathSub p} the honest sibling path, and \lean{t.get},
\lean{t.set}, \lean{t.path} their leaf-level specialisations.

\begin{lstlisting}
theorem fold_pathSub (t : Tree d) (p : Pos d h) :
    fold (t.getSub p).root (t.pathSub p) = t.root

theorem fold_setSub (t : Tree d) (p : Pos d h) (s : Tree h) :
    fold s.root (t.pathSub p) = (t.setSub p s).root

theorem root_inj (t₁ t₂ : Tree d) : t₁.root = t₂.root → t₁ = t₂

theorem fold_sound (t : Tree d) (m : Node H) (π : List (Bool × Node H)) :
    fold m π = t.root →
    ∃ (h : Nat) (p : Pos d h), m = (t.getSub p).root ∧ π = t.pathSub p
\end{lstlisting}

\lean{fold\_sound} is Theorem~\ref{thm:binding}. Its proof is the only
place the hash axioms are used: a non-empty path claims an internal node,
so the tree is not a tip by tag separation, and injectivity of
\lean{hNode} equates both children's hashes and sums with the claimed
ones. Note that the conclusion equates \emph{nodes}, so the sum the
prover presented is the true sum of the subtree.

A \lean{Step} is a path, a height, an old node and a new node; it is
valid for roots $r, r'$ when the old node folds to $r$ and the new node to
$r'$, and well formed for depth $d$ when the path has $d - h$ entries.
\lean{Steps} chains steps, each against the root the previous one produced.

\begin{lstlisting}
theorem Step.sound (t : Tree d) (s : Step H) (r' : Node H)
    (hwf : s.WF d) (hv : s.valid t.root r') :
    ∃ p : Pos d s.height, s.old = (t.getSub p).root ∧
      ∀ (u : Tree s.height), u.root = s.new → r' = (t.setSub p u).root

theorem Steps.sum (h : Steps r ss r') :
    r'.sum + oldSum ss = r.sum + newSum ss

theorem Steps.sound (t : Tree d) (h : Steps t.root ss r')
    (hwf : ∀ s ∈ ss, s.WF d) (hg : ∀ s ∈ ss, s.Grounded) :
    ∃ t' : Tree d, t'.root = r'
\end{lstlisting}

\lean{Steps.sum} is Theorem~\ref{thm:cons} and needs no hash assumption.
\lean{Steps.sound} is Theorem~\ref{thm:ground}; \lean{s.Grounded} says
the new node is the root of some tree of the step's height, which
leaf-level steps satisfy by construction (\lean{grounded\_of\_leaf}) and
which a \tr{JOIN} graft must be given.

The transitions are structures bundling the steps with the output rule
and an opaque authorisation. Their theorems:

\begin{lstlisting}
theorem payable (t : Tree d) (p : Pos d 0) (l : Leaf) (hl : t.get p = some l) :
    (leafStep (t.path p) (some l) none).valid t.root (t.set p none).root ∧
    (t.set p none).root.sum + l.amount = t.root.sum ∧
    ∀ q, q ≠ p → (t.set p none).get q = t.get q

theorem ExitK.no_inflation (h : ExitK exact r blocks r' out0) :
    out0 + payoutSum blocks ≤ r.sum

theorem ExitK.exact_conservation (h : ExitK true r blocks r' out0) :
    out0 + payoutSum blocks = r.sum

theorem join_sound (t : Tree d) (hwf : path.length + h = d)
    (hj : Join auth h t.root path new r' out0) :
    ∃ p : Pos d h, t.getSub p = Tree.empty h ∧
      out0 = t.root.sum + new.sum ∧
      ∀ u : Tree h, u.root = new → r' = (t.setSub p u).root

theorem claim_sound (t : Tree d) (hwf₁ : srcPath.length = d)
    (hwf₂ : claimPath.length = d)
    (hc : Claim auth t.root srcPath claimPath src claim r' out0) :
    out0 = t.root.sum ∧
    ∃ (p₁ p₂ : Pos d 0), t.get p₁ = some src ∧
      (t.set p₁ (claimedSource src claim.amount)).get p₂ = none ∧
      r' = ((t.set p₁ (claimedSource src claim.amount))
              .set p₂ (some claim)).root
\end{lstlisting}

\lean{payable} is Theorem~\ref{thm:pay}. \lean{join\_sound} uses
\lean{root\_inj} (via \lean{eq\_empty\_of\_root}): since \tr{JOIN}$_h$
hard-codes $\Emp_h$ as the old node, the slot it fills is provably the
empty tree, and the theorem gives the successor root for every tree whose
root is the grafted node. \lean{claim\_sound} exposes the drained-source
case: the claim position must be empty in the tree \emph{after} the source
step, which is why a source reduced to zero may host the claim.

\subsection{Assumption Ledger}

Every hypothesis outside Lean's logic, what it stands for in the
deployment, and which theorems consume it.

\begin{center}\scriptsize
\begin{tabular}{@{}p{3.6cm}p{4.2cm}p{3.6cm}@{}}
\toprule
hypothesis & stands for & used by \\
\midrule
\lean{hLeaf\_inj} & SHA-256 collision resistance + \lean{serializeLeaf\_inj} & \lean{leafNode\_inj}, \lean{root\_inj} \\
\lean{hNode\_inj} & SHA-256 collision resistance + \lean{serializeNode\_inj} & \lean{root\_inj}, \lean{fold\_sound} \\
\lean{hNode\_ne\_hLeaf}, \lean{hNode\_ne\_empty}, \lean{hLeaf\_ne\_empty} & distinct tags (random-oracle heuristic) & \lean{leafNode\_inj}, \lean{fold\_sound} \\
\lean{treeSpk\_inj} & DL on secp256k1 + CR of $\Hstate$, TapTweak & \lean{witness\_root\_binds} \\
real-tree invariant & creation check + grounded grafts (off-chain) & every theorem stated for a \lean{Tree} \\
\midrule
(none) & & \lean{Steps.sum}, \lean{ExitK.no\_inflation}, encodings \\
\bottomrule
\end{tabular}
\end{center}

The conservation results use no hypothesis at all, which is the formal
content of ``the covenant conserves sums whatever it is shown''. The
collision-resistance assumption on the TapSighash and challenge hashes
that the MIN tier's constructive introspection adds (\S\ref{sec:min}) is
outside the model, since the model does not include signature checking.

\subsection{Modules and Effort}

\begin{center}\scriptsize
\begin{tabular}{@{}lrp{7.5cm}@{}}
\toprule
module & lines & contents \\
\midrule
\lean{Hash.lean} & 59 & the \lean{HashModel} class \\
\lean{Leaf.lean} & 96 & \lean{Leaf}, \lean{Leaf.WF}, \lean{Node}, \lean{leafNode}, \lean{leafNode\_inj} \\
\lean{Tree.lean} & 335 & trees, positions, paths, \lean{fold}, all tree theorems \\
\lean{Transitions.lean} & 458 & \lean{Step}, \lean{Steps}, the transitions and their theorems \\
\lean{Binding.lean} & 60 & \lean{BindingModel} and the two binding theorems \\
\lean{Encoding.lean} & 238 & encodings, serialisation injectivity, \lean{SEQ\_FINAL} \\
\lean{Axioms.lean} & 32 & the axiom audit \\
\bottomrule
\end{tabular}
\end{center}

The only proofs that took real effort were \lean{fold\_sound}, where the
two-index positions and the domain-separation case for a tip had to be
found, and \lean{Steps.sound}, where the height of the grounding tree
must be reconciled with the height the soundness lemma returns, which is
why \lean{Step} carries its height explicitly. A verified interpreter for
the tapscript subset would need a model of the stack machine, of
BIP-441 numerics, of \op{TX}'s serialisations and of the BIP-341 message,
and a proof that each generated script's execution relation equals the
corresponding \lean{Step} relation; we estimate it at several thousand
lines and regard it as the natural sequel.

\subsection{Binding to the Output}

\lean{Binding.lean} records the one cryptographic step outside the hash
model. A structure \lean{BindingModel} carries an abstract
\lean{treeSpk : H → TapRoot → Spk} with the hypothesis that it is
injective in both arguments, and two one-line theorems conclude that a
witness $(r', m')$ passing the scriptPubKey check equals the committed
$(r, m)$, and that the successor output commits to the computed root under
the same $m$. In the deployment the map is
$(r, m) \mapsto x\bigl(P_r + H_{\mathsf{TapTweak}}(x(P_r) \,\|\, m)\cdot G\bigr)$
with $P_r = \mathsf{lift\_x}(x(\NUMS + \Hstate(r) G))$, and the two
x-projections identify a point with its negation. A collision
$(r, m) \ne (r', m')$ yields one of: $x(P_r) = x(P_{r'})$ with $r \ne r'$,
which is a collision in $\Hstate$ or the relation
$\NUMS + tG = -(\NUMS + t'G)$, solvable for the discrete logarithm of
$\NUMS$; $x(P_r) = x(P_{r'})$ with $m \ne m'$, a collision in
$H_{\mathsf{TapTweak}}$; or distinct internal keys with $Q = \pm Q'$,
again a linear relation in $\NUMS$ and $G$ solvable for the discrete
logarithm of $\NUMS$. Making the hypothesis explicit is the point: it is
what a future script change must not break.

\subsection{Encodings}

\lean{Encoding.lean} proves the byte-level facts the design leans on, over
natural-number bytes so that every proof stays within linear arithmetic:
fixed-width little-endian encodings are injective on their range;
\lean{minimalLE (n + 2\^{}(8k)) = leBytes k n ++ [1]} for $n < 2^{8k}$,
which with $k = 8$ is $\encs$ and justifies producing it with one
addition; the 116-byte leaf and the 82-byte node preimage are injective;
and the \lean{SEQ\_FINAL} facts of Lemma~\ref{lem:enc}, computed by
\lean{decide} on a bounded-recursion model of \texttt{CScriptNum}.

\subsection{What Is Not Proved}

The tapscript bytecode is outside the model. That the scripts of
\S\ref{sec:scripts} implement the relations above is established by
execution: every script of every tier runs in the simulator against a
reference tree with positive and negative cases, and the Go generators are
compared byte for byte with the Python output. A verified interpreter for
the tapscript subset (BIP-441 numerics, \op{TX}, \op{CHECKSIGFROMSTACK})
would close this gap and is a project of a different size. Signatures,
lock times and the transaction structure beyond amounts are
authorisation-side facts and are deliberately opaque. Concurrency,
mempool policy, reorganisations and the off-chain protocol are not
theorems about hash trees and belong to a model checker; a design for a
P~\cite{plang} model is in the artefact.
